\Appendix

\chapter*{附录1}

\begin{center}
  \large 问卷调查
\end{center}

\includegraphics[height = 0.85\textheight]{appendix-1}

\includegraphics[height = 0.9\textheight]{appendix-2}

\includegraphics[width = 0.8\textwidth]{appendix-3}
% 性别: $\square$男 \quad $\square$女

% \begin{enumerate}[1.]
%   \item 我能够根据问题中的量关系写出数学表达式
%   \par\begin{tblr}{ccccc}
%     1非常同意 & 2同意 & 3不确定 & 4不同意 &  5非常不同意
%   \end{tblr}
%   \item 我能对熟悉的数学应用与建模的结果作出说明或解释
%   \par\begin{tblr}{ccccc}
%     1非常同意 & 2同意 & 3不确定 & 4不同意 &  5非常不同意
%   \end{tblr}
%   \item 我能清楚记忆常见的数学公式和法则
%   \par\begin{tblr}{ccccc}
%     1非常同意 & 2同意 & 3不确定 & 4不同意 &  5非常不同意
%   \end{tblr}
%   \item 在学习运算法则时，我会思考该法则的意义和作用
%   \par\begin{tblr}{ccccc}
%     1非常同意 & 2同意 & 3不确定 & 4不同意 &  5非常不同意
%   \end{tblr}
%   \item 在完成数学运算后，我会思考其他的运算方法
%   \par\begin{tblr}{ccccc}
%     1非常同意 & 2同意 & 3不确定 & 4不同意 &  5非常不同意
%   \end{tblr}
%   \item 我能够运用运算验证简单的数学结论
%   \par\begin{tblr}{ccccc}
%     1非常同意 & 2同意 & 3不确定 & 4不同意 &  5非常不同意
%   \end{tblr}
%   \item 在进行数学运算时,我能合理选择公式，并灵活运用
%   \par\begin{tblr}{ccccc}
%     1非常同意 & 2同意 & 3不确定 & 4不同意 &  5非常不同意
%   \end{tblr}
%   \item 我在拿到数学题目后会认真审题
%   \par\begin{tblr}{ccccc}
%     1非常同意 & 2同意 & 3不确定 & 4不同意 &  5非常不同意
%   \end{tblr}
%   \item 我能规范地书写运算过程
%   \par\begin{tblr}{ccccc}
%     1非常同意 & 2同意 & 3不确定 & 4不同意 &  5非常不同意
%   \end{tblr}
%   \item 我能指岀某一数学问题是向量运算问题
%   \par\begin{tblr}{ccccc}
%     1非常同意 & 2同意 & 3不确定 & 4不同意 &  5非常不同意
%   \end{tblr}
%   \item 我能够选择相应的运算法则来进行计算
%   \par\begin{tblr}{ccccc}
%     1非常同意 & 2同意 & 3不确定 & 4不同意 &  5非常不同意
%   \end{tblr}
%   \item 我能够将运算方法运用到相关情境中解决问题
%   \par\begin{tblr}{ccccc}
%     1非常同意 & 2同意 & 3不确定 & 4不同意 &  5非常不同意
%   \end{tblr}
%   \item 我能够和同学交流思路，探讨更为合理简洁 的运算方法
%   \par\begin{tblr}{ccccc}
%     1非常同意 & 2同意 & 3不确定 & 4不同意 &  5非常不同意
%   \end{tblr}
%   \item 我平时在做完数学运算后，会再检查一下
%   \par\begin{tblr}{ccccc}
%     1非常同意 & 2同意 & 3不确定 & 4不同意 &  5非常不同意
%   \end{tblr}
%   \item 当运算进行不下去时，我会改变算法或者寻求别的方法
%   \par\begin{tblr}{ccccc}
%     1非常同意 & 2同意 & 3不确定 & 4不同意 &  5非常不同意
%   \end{tblr}
%   \item 在解题时，我更倾向于先思考一会而不是直接动笔
%   \par\begin{tblr}{ccccc}
%     1非常同意 & 2同意 & 3不确定 & 4不同意 &  5非常不同意
%   \end{tblr}
%   \item 我能够利用分类讨论解决复杂的代数运算问题
%   \par\begin{tblr}{ccccc}
%     1非常同意 & 2同意 & 3不确定 & 4不同意 &  5非常不同意
%   \end{tblr}
%   \item 我会与同学交流某道题的运算结果
%   \par\begin{tblr}{ccccc}
%     1非常同意 & 2同意 & 3不确定 & 4不同意 &  5非常不同意
%   \end{tblr}
%   \item 在综合运算问题中，我会因运算不规范而失分
%   \par\begin{tblr}{ccccc}
%     1非常同意 & 2同意 & 3不确定 & 4不同意 &  5非常不同意
%   \end{tblr}
%   \item 在解决有关运算问题时，我能联想类似的运算方法
%   \par\begin{tblr}{ccccc}
%     1非常同意 & 2同意 & 3不确定 & 4不同意 &  5非常不同意
%   \end{tblr}
%   \item 对运算中的错误，我会寻找原因，并进行反思
%   \par\begin{tblr}{ccccc}
%     1非常同意 & 2同意 & 3不确定 & 4不同意 &  5非常不同意
%   \end{tblr}
%   \item 我会自己尝试总结-类问题的运算技巧或方法
%   \par\begin{tblr}{ccccc}
%     1非常同意 & 2同意 & 3不确定 & 4不同意 &  5非常不同意
%   \end{tblr}
%   \item 我能由数思形思考代数运算问题
%   \par\begin{tblr}{ccccc}
%     1非常同意 & 2同意 & 3不确定 & 4不同意 &  5非常不同意
%   \end{tblr}
%   \item 我能够针对不同类型问题的特点，选择合理简洁的运算途径
%   \par\begin{tblr}{ccccc}
%     1非常同意 & 2同意 & 3不确定 & 4不同意 &  5非常不同意
%   \end{tblr}
% \end{enumerate}



\chapter*{附录2}

\begin{center}
  \large
  访谈题目
\end{center}
% \begin{enumerate}[\chinese*、]
%   \item 根据您多年的教学经验，您觉得现在的西藏高中生在数学学习中存在哪些问题？

%   \item 您认为现阶段西藏高中生的数学运算能力存在哪些问题?
  
%   \item 您认为教师日常教学中重视学生运算能力的培养吗？
  
%   \item 您认为应该怎样有效地培养西藏高中生数学运算能力？
% \end{enumerate}

\includegraphics[width = \textwidth]{appendix-2-1}
